\section{A width-independent tanh bound without robustness}
\label{sec:productnorm}

We first establish the width-independent function-class bound used to
control the low-energy signal extracted from a posterior draw. The proof
accounts for the fact that tanh is not positively homogeneous.

The normalized-activation and exponential-moment method is closely related
to Nakhleh--Nowak's treatment of deep neural variation spaces
\cite{nakhleh2026}, which explicitly uses $\tanh(st)/s$ and slope
discretization, and to the exponential peeling argument of
Golowich--Rakhlin--Shamir~\cite{golowich2018}. We give a self-contained
product-Frobenius version for the present Gaussian-energy conventions;
we do not claim the underlying proof technique is new.

\begin{theorem}[Product-Frobenius bound for tanh]\label{thm:tanh-product}
Let $\mathcal F_{L,P}$ contain all scalar-output tanh networks with $L-1$
hidden layers, arbitrary finite widths, a linear readout, and product of
layer Frobenius norms at most $P$. Biases are absent except through an
augmented input. On every sample $x_1,\ldots,x_n$ with $\|x_i\|\le R$,
\[
 \widehat{\mathfrak R}_n(\mathcal F_{L,P})
 \le \frac{PR}{\sqrt n}
 \left[2+\sqrt{2(L-1)\log\bigl(3+n(L-1)^2\bigr)}\right].
\]
The bound is independent of every hidden width and of the individual layer
norms apart from their product.
\end{theorem}

\begin{proof}
If $R=0$, every sampled input and every network output on the sample is
zero, so the bound is immediate. Assume $R>0$ below.
Define $\phi_s(u)=\tanh(su)/s$ for $s>0$, $\phi_0(u)=u$, and
$\phi_\infty(u)=0$. These functions are odd, $1$-Lipschitz, and zero at zero.
For finite $s,t\ge0$ and $r\ge0$,
\[
 \phi_s(ru)=r\phi_{sr}(u),\qquad
 |\phi_s(u)-\phi_t(u)|\le u^2|s-t|,
 \qquad |\phi_s(u)|\le1/s\quad(s>0).
\]
The derivative bound follows by differentiating
$\phi_s(u)=u\int_0^1\operatorname{sech}^2(suv)\,dv$.

Let $\mathcal G_0$ contain the identity feature map. Inductively allow
arbitrary finite-dimensional features
\[
 g_j(x)=\phi_{s_j}(a_j^T h(x)),\quad
 h\in\mathcal G_{\ell-1},\quad \sum_j\|a_j\|^2\le1,
 \quad s_j\in[0,\infty].
\]
All these feature vectors have norm at most $R$ on the sample. Put
$B_\ell(\epsilon)=\sup_{g\in\mathcal G_\ell}
 \|\sum_i\epsilon_i g(x_i)\|$ for independent Rademacher signs.
Writing $a_j=r_j u_j$ and using the scaling identity shows
\[
 B_\ell\le
 \sup_{h,\,\|u\|\le1,\,s\in[0,\infty]}
 \left|\sum_i\epsilon_i\phi_s(u^T h(x_i))\right|.
\]
For any $\eta>0$, uniformly on $[-R,R]$ the slope family admits an
$\eta/n$ approximation using at most
$J=3+n^2R^2/\eta^2$ choices: truncate $s>n/\eta$ to the zero function,
and round slopes below that threshold down at mesh $\eta/(nR^2)$.
The error in the Rademacher sum is at most $\eta$.

We use the convex contraction principle
\[
 \mathbb E G\!\left(\sup_{t\in T}\sum_i\epsilon_i\psi_i(t_i)\right)
 \le
 \mathbb E G\!\left(\sup_{t\in T}\sum_i\epsilon_i t_i\right)
\]
for convex nondecreasing $G$ and $1$-Lipschitz $\psi_i$ vanishing at zero.
For completeness, the single-coordinate proof compares
$X_\pm=\sup_t(a_t\pm\psi(t_k))$ and
$Y_\pm=\sup_t(a_t\pm t_k)$. One has
$\max(Y_+,Y_-)\ge\max(X_+,X_-)$ by $|\psi(t)|\le|t|$.
An approximately maximizing pair for $X_+,X_-$ also gives
$Y_++Y_-\ge X_++X_-$, by the Lipschitz property.
Convex monotonicity therefore gives
$G(Y_+)+G(Y_-)\ge G(X_+)+G(X_-)$; iterate over coordinates.
Bounded separable classes follow by approximation from finite classes.
In our application, $u$ ranges symmetrically and $\phi_s$ is odd, so the
absolute supremum is the same as the one-sided supremum.

Apply this principle with $G(t)=e^{\lambda t}$ to each grid slope, then
bound the exponential of a maximum by the sum of the exponentials:
\[
 \mathbb E e^{\lambda B_\ell}
 \le e^{\lambda\eta}J\,\mathbb E e^{\lambda B_{\ell-1}}.
\]
Furthermore, $\mathbb E B_0\le R\sqrt n$, and bounded differences gives
$\mathbb E e^{\lambda B_0}
 \le e^{\lambda R\sqrt n+\lambda^2nR^2/2}$.
Writing $H=L-1$, iteration and Jensen yield
\[
 \mathbb E B_H\le H\eta+R\sqrt n+
             \frac{H\log J}{\lambda}+\frac{\lambda nR^2}{2}.
\]
Take $\eta=R\sqrt n/H$ and optimize $\lambda$ to obtain
$\mathbb E B_H\le R\sqrt n[2+\sqrt{2H\log(3+nH^2)}]$.

Finally, for a particular original network, write
$M_\ell=\|W^\ell\|_F$ and $S_\ell=\prod_{j\le\ell}M_j$.
Dividing its $\ell$th hidden vector by $S_\ell$ changes its recursion to
$\phi_{S_\ell}((W^\ell/M_\ell)g^{\ell-1})$.
Thus its normalized hidden features belong to the enlarged classes above,
and its scalar output is at most $P$ times a norm-one linear readout of
$\mathcal G_H$. Zero-norm layers give the zero function and require no
normalization. This proves the theorem. All suprema are measurable by
countable approximation of finite-dimensional parameter spaces and then
union over widths.
\end{proof}

For the Gaussian energy in our model, the product of the $L$ weighted
squared layer norms is $dP^2$. Consequently
\[
 P\le d^{-1/2}(\kappa_N/L)^{L/2}
 \le d^{-1/2}e^{B/2}\qquad\text{if }\kappa_N\le L+B.
\]
Applying the theorem to the ramp loss
$\ell(yz)=\min\{1,\max\{0,1-2yz\}\}$ gives the following explicit bound.

\begin{corollary}[All widths at bounded excess energy]\label{cor:excess-risk}
Fix $B\ge0$. With probability at least $1-\delta$, every sample-consistent
depth-$L$ network, at any finite width, with training margins at least $1/2$
and energy $\kappa_N\le L+B$ has sign error at most
\[
 \frac{4\sqrt2 e^{B/2}}{\sqrt n}
 \left[2+\sqrt{2(L-1)\log\bigl(3+n(L-1)^2\bigr)}\right]
 +\sqrt{\frac{\log(1/\delta)}{2n}}.
\]
No internal robustness assumption is imposed.
\end{corollary}

\begin{proof}
Augmented Boolean inputs have norm $\sqrt{d+1}$, so $PR\le\sqrt2e^{B/2}$.
The ramp loss is bounded by one and is $2$-Lipschitz. Symmetrization,
scalar contraction, and bounded differences give population loss at most
empirical loss plus four times the network Rademacher bound and the
stated confidence term. Empirical loss is zero and population loss
upper-bounds sign error.
\end{proof}

\begin{lemma}[A scalar transmission construction]\label{lem:scalar-transmission}
For every real $D>0$ and integer $m\ge2$, there are positive scalar weights
$\alpha,b_2,\ldots,b_{m-1},c$ such that
\[
 r_1=\tanh\alpha,\quad r_s=\tanh(b_s r_{s-1}),\quad
 cr_{m-1}=1/2,
\]
and, writing $D'=\max\{D,4\}$,
\[
 D\alpha^2+\sum_{s=2}^{m-1}b_s^2+c^2
 \le(m+12)(D'/4)^{1/m},\qquad c^2=(D'/4)^{1/m}.
\]
\end{lemma}

\begin{proof}
Set $q=(D'/4)^{1/(2m)}$, $k_s=\min\{s,m-s\}$, and
$r_s=q^s/(\sqrt{D'}k_s)$. Then $0<r_s\le1/(2qk_s)\le1/2$ and
$r_{m-1}=1/(2q)$. Choose
$\alpha=\operatorname{atanh}r_1$,
$b_s=\operatorname{atanh}(r_s)/r_{s-1}$ and $c=q$.
For $0<u\le1/2$ the power series gives
$(\operatorname{atanh}(u)/u)^2\le1+u^2$.
Writing $R_s=k_{s-1}/k_s$, the reciprocal pairs
$k/(k+1),(k+1)/k$ satisfy
\[
 \left(\frac{k+1}{k}\right)^2+
 \left(\frac{k}{k+1}\right)^2-2\le4/k^2.
\]
Consequently $\sum_{s=2}^{m-1}R_s^2\le m-2+7$,
$R_s\le2$, and $\sum_{s=1}^{m-1}k_s^{-2}<4$.
Using $q^{2m}=D'/4$ therefore gives
\[
 D\alpha^2\le q^2+1/4,\qquad
 \sum_{s=2}^{m-1}b_s^2\le q^2(m-2+7)+4,\qquad c^2=q^2.
\]
Their sum is at most $(m+12)q^2$. Empty internal sums at $m=2$
satisfy the same estimates.
\end{proof}

